\documentclass[11pt]{article}
\usepackage[margin=1in]{geometry}
\usepackage{amsmath,amssymb,amsthm}
\usepackage[T1]{fontenc}
\usepackage{hyperref}
\hypersetup{colorlinks=true,linkcolor=blue,urlcolor=blue}
\newtheorem{theorem}{Theorem}
\newtheorem{lemma}[theorem]{Lemma}
\newtheorem{corollary}[theorem]{Corollary}
\newtheorem{remark}[theorem]{Remark}
\newcommand{\R}{\mathbb R}
\newcommand{\Q}{\mathbb Q}
\newcommand{\PC}{\operatorname{PC}}
\newcommand{\relint}{\operatorname{relint}}
\newcommand{\vol}{\operatorname{vol}}
\title{Positive balancing, compact support, and a polyhedral maximum principle\\\large A supporting scope calculation for P05.1 and P05.2}
\author{Research draft}
\date{7 October 2026}
\begin{document}
\maketitle
\begin{abstract}
We give elementary proofs that a nonempty finite positive balanced rational cycle of positive dimension in one real vector space has noncompact support, and that the convex cone generated by a positive balanced fan is its linear span. The latter yields a strong maximum principle for continuous functions satisfying the rational-endpoint polyhedral Jensen inequalities. On compact connected spaces with compatible local balanced charts, these functions are constant. These statements clarify the distinction between a globally embedded compact support, a finite bounded skeleton, a toric compactification, and an abstract compact space with local charts. No regularity theorem or higher-dimensional contact classification is asserted.
\end{abstract}

\section{Exact geometric conventions}
Let $N$ be a lattice, $V=N\otimes\R$, and fix an auxiliary Euclidean inner product on $V$. A finite rational polyhedral complex $\Pi$ is pure of dimension $d$. Its tangent lattices are saturated:
\[N_\sigma=N\cap\operatorname{span}_{\R}(\sigma-\sigma).
\]
For a facet $\tau\prec\sigma$, choose an inward integral vector $n_{\sigma/\tau}$ whose class is primitive in $N_\sigma/N_\tau$. A strictly positive balanced weight $w$ satisfies
\begin{equation}\label{balance}
\sum_{\sigma\succ\tau}w_\sigma n_{\sigma/\tau}\in N_{\tau,\R}.
\end{equation}
Weights may be positive real numbers; rational weights in the fixed BMPS conventions are a special case. The definitions of finite globally embedded space, saturated lattices, inward normals and strictly positive top cycles agree with \cite[Definitions 2.1--2.4, 2.18--2.27]{BMPS}. Only definitions are imported. No source regularization or convexity theorem is a grant.

For a rational affine face $\sigma$, let $a_\sigma>0$ be the Euclidean covolume of $N_\sigma$ in its tangent space, and put $a_\tau=1$ for a zero-dimensional face. Set
\[c_\sigma=w_\sigma a_\sigma.
\]
These auxiliary weights serve an ordinary Euclidean divergence calculation. They are not a replacement for the raw polyhedral intersection weights or a Monge--Amp\`ere measure.

\begin{lemma}\label{height}
Let $\pi_\tau$ denote orthogonal projection onto $N_{\tau,\R}^{\perp}$. If $\nu_{\sigma/\tau}$ is the inward Euclidean unit conormal in the tangent space of $\sigma$, then
\[
 \pi_\tau n_{\sigma/\tau}
   =\frac{a_\sigma}{a_\tau}\nu_{\sigma/\tau},
 \qquad
 \sum_{\sigma\succ\tau}c_\sigma\nu_{\sigma/\tau}=0.
\]
\end{lemma}
\begin{proof}
The saturated rank-$d-1$ sublattice $N_\tau\subset N_\sigma$ has torsion-free rank-one quotient. A basis of $N_\tau$ followed by $n_{\sigma/\tau}$ is therefore a basis of $N_\sigma$. Its parallelotope volume is the base volume $a_\tau$ times the perpendicular height $|\pi_\tau n_{\sigma/\tau}|$. The inward sign gives the first formula. Project \eqref{balance} orthogonally and multiply by $a_\tau$ to obtain the second. Tangential changes of a primitive lift have no effect.
\end{proof}

\begin{lemma}\label{stationary}
For any compactly supported continuously differentiable ambient vector field $Y$,
\begin{equation}\label{stationarity}
 \sum_{\sigma\in\Pi(d)}c_\sigma
       \int_\sigma\operatorname{div}_{\sigma}Y\,d\vol_d=0.
\end{equation}
Here the divergence is the trace of the derivative on the affine tangent space of $\sigma$.
\end{lemma}
\begin{proof}
Apply the ordinary divergence formula to each polyhedron. For unbounded cells, truncate outside the support of $Y$; the added boundary gives zero. Each genuine facet $\tau$ contributes the integral of $Y$ paired with the outward unit conormal, which is $-\nu_{\sigma/\tau}$. Its total coefficient is zero by Lemma~\ref{height}. Lower-dimensional faces have zero facet measure. For $d=1$, the same calculation is the fundamental theorem of calculus on each edge, and the zero-dimensional facet integrals are endpoint values. This proves \eqref{stationarity}. The divergence formula itself follows by triangulation and the one-variable fundamental theorem, so it introduces no polyhedral intersection hypothesis.
\end{proof}

\section{A compact-support obstruction}
\begin{theorem}\label{noncompact}
If $\Pi$ is nonempty, finite, pure of dimension $d\geq1$, and admits strictly positive balanced top weights, then its support $|\Pi|\subset V$ is noncompact.
\end{theorem}
\begin{proof}
If the support were compact, every face would be bounded. Choose a compactly supported vector field equal to $Y(x)=x$ on a neighborhood of the support. On every top face its tangential divergence is $d$. Equation~\eqref{stationarity} would give
\[0=d\sum_{\sigma\in\Pi(d)}c_\sigma\vol_d(\sigma)>0.
\]
The sum is positive because the weights and covolumes are positive and each nonempty top polyhedron has positive relative volume. This is a contradiction.
\end{proof}

\begin{remark}\label{scope}
This theorem concerns the original support in a single $V$. It applies directly to the globally embedded finite spaces of BMPS Definition 2.4 equipped with Definition 2.27's strictly positive cycle. A union of bounded faces inside a noncompact balanced space need not be balanced at its boundary. A toric compactification has additional strata and lies outside this single-vector-space hypothesis. An abstract compact space with local integral-affine charts also lies outside that hypothesis. Consequently no assertion about these other categories follows by naming them compact supports. In particular the theorem supplies no higher-dimensional real envelope/contact classification.
\end{remark}

\section{The cone generated by a positive fan}
\begin{theorem}\label{fan}
Let $\Sigma$ be a nonempty finite rational fan, allowing common lineality, pure of dimension $d\geq1$, with strictly positive balanced weights. If $L$ is the linear span of its support, then
\[\operatorname{cone}(|\Sigma|)=L.
\]
Equivalently, no linear form is nonnegative on the entire support and positive somewhere on it.
\end{theorem}
\begin{proof}
Choose a smooth nonincreasing radial cutoff $\chi(r)$, equal to one near zero, zero for large $r$, and decreasing on a nonempty interval. For any constant vector $b\in V$, apply Lemma~\ref{stationary} to $Y(x)=b\chi(|x|)$. For $x$ in a top cone, its radial vector belongs to that cone's tangent space, and
\[\operatorname{div}_\sigma Y(x)
      =\chi'(|x|)\frac{\langle b,x\rangle}{|x|}.
\]
Polar integration on the $d$-dimensional tangent space gives
\[0=\left(\int_0^\infty\chi'(r)r^{d-1}\,dr\right)
 \sum_{\sigma\in\Sigma(d)}c_\sigma
     \int_{\sigma\cap S(V)}\langle b,\theta\rangle\,dS_{d-1}(\theta).
\]
The radial factor is finite and strictly negative, including when $d=1$, for which the spherical integral is a sum over the directions. Thus the weighted angular vector has barycenter zero.

If a linear form $\ell$ is nonnegative on the support and positive at some point, a top cone containing that point has a relatively open angular subset where $\ell>0$. Its angular measure is positive. All other angular integrals are nonnegative, and all $c_\sigma$ are positive, contradicting the zero barycenter. Hence every such $\ell$ vanishes on the support.

The cone generated by the support is a closed polyhedral cone: each of the finitely many rational cones is finitely generated, with a basis and its negatives added for any lineality. If that cone were a proper subset of $L$, finite-dimensional separation would give a linear form nonnegative on it and nonzero on $L$. It is positive somewhere on the support because the support spans $L$, a contradiction. For completeness, separation of a point from a closed convex cone follows by taking its nearest point: the vector from the outside point to the nearest point pairs nonnegatively with the cone and negatively with the outside point. This proves the equality.
\end{proof}

\begin{corollary}\label{relation}
Choose any finite set of nonzero rational vectors $v_1,\ldots,v_m$ generating all cones of $\Sigma$, including both signs for lineality. There exist positive rational numbers $\lambda_i$ with
\[\sum_{i=1}^m\lambda_i v_i=0.
\]
\end{corollary}
\begin{proof}
Theorem~\ref{fan} puts each $-v_i$ in the cone generated by the $v_j$. Its nonnegative representation can be chosen rational: eliminate dependent active columns until the remaining columns are linearly independent, then solve that rational linear system. Adding the resulting relations $v_i+\sum_j b_{ij}v_j=0$ gives the desired rational relation, with every coefficient at least one.
\end{proof}

\section{A continuous polyhedral maximum principle}
On a finite rational structure, a rational-endpoint polyhedral convex combination is an equality
\[x=\sum_i\alpha_i x_i,\qquad \alpha_i\geq0,\quad\sum_i\alpha_i=1,
\]
where the endpoints $x_i$ are rational and belong to faces $\sigma_i$ containing a face $\tau$ with $x\in\tau$. A continuous function $f$ is PC here if it satisfies the associated Jensen inequality for every such combination and every rational subdivision. These are BMPS Definitions 3.2--3.3 with continuity made an explicit hypothesis. We do not import its cited automatic-continuity theorem.

\begin{remark}\label{rational-pathology}
The printed rational endpoint definition does not ensure continuity. On $\R$ with its usual rational polyhedral structure, put $f(x)=0$ for $x\in\Q$ and $f(x)=-1$ otherwise. Every permitted endpoint has value zero, and every center has value at most zero, so every permitted PCC Jensen inequality holds. Nevertheless $f$ is discontinuous at every point. The source's sentence after Definition 3.3 asserting automatic continuity therefore requires an additional hypothesis or a change of endpoint convention. This example changes neither the definition tested nor the continuous class of the following theorem. It makes no assertion that an envelope of continuous obstacles is discontinuous.
\end{remark}

\begin{theorem}\label{maximum}
On a positive balanced finite rational polyhedral space, a continuous PC function attaining a global maximum is constant on the connected component of each maximum point. The same assertion holds for a space with local rational polyhedral charts and compatible positive balanced weights, provided the function satisfies these local Jensen inequalities. In particular every continuous PC function on a compact connected space of the latter kind is constant.
\end{theorem}
\begin{proof}
In dimension zero every connected component is a point, so the conclusion is immediate. In positive dimension first consider a rational maximum point $x\in\relint(\tau)$. Its local star consists of the cones $C_{\sigma/\tau}$ in the ambient tangent vector space, with their inherited weights. They are balanced: faces containing $\tau$ and their inward primitive normal classes are exactly the corresponding faces and normal classes of the original structure. Their tangent lattices are unchanged. Apply Corollary~\ref{relation} to nonzero rational generators $v_i$ of these cones, including both directions along $\tau$. Choose a rational $\varepsilon>0$ small enough that $z_i=x+\varepsilon v_i$ belongs to an incident face $\sigma_i\supset\tau$.

Let $y$ be any rational point of an incident face $\sigma\supset\tau$. Its displacement is in $C_{\sigma/\tau}$, so write $y-x=\sum_i a_i v_i$ with $a_i\geq0$. Choose $t$ large enough that $t\lambda_i-a_i\geq0$ for every $i$. Then
\[(y-x)+\sum_i\frac{t\lambda_i-a_i}{\varepsilon}(z_i-x)=0.
\]
After normalizing the coefficients, this expresses $x$ as a PCC of $y$ and the $z_i$, with a strictly positive coefficient at $y$. All endpoint values are at most the global maximum $M=f(x)$. Jensen's inequality forces $f(y)=M$. Rational points are dense in the incident faces, so continuity implies that the entire closed star of $\tau$ has value $M$.

For a possibly irrational maximum point $p\in\relint(\tau)$, the restriction of $f$ to $\tau$ is ordinarily convex. Indeed rational-endpoint Jensen inequalities extend to arbitrary endpoints by continuity. A convex function with a maximum at a relative interior point is constant on the whole face: for any $q\in\tau$, express $p$ as a convex combination with positive coefficient at $q$ and a second point in $\tau$ on the other side of $p$, and use the global upper bound $M$. Thus a rational point $x\in\relint(\tau)$ also has value $M$. The preceding argument makes every incident face maximal, giving a neighborhood of $p$ in the maximum set.

The maximum set is therefore open and, by continuity, closed. Its intersection with a connected component containing a maximum point is that entire component. In local charts, the same argument uses only a sufficiently small finite star and the global upper bound; every endpoint used lies in the chart. Compactness supplies a maximum on each compact connected component, proving the last assertion.
\end{proof}

\section{Consequences and limits for the envelope problems}
The compact globally embedded positive-dimensional branch is excluded by Theorem~\ref{noncompact}, independently of the reference or coefficient convention. This does not exclude abstract compact balanced spaces, and Theorem~\ref{maximum} gives a separate, explicitly continuous statement there.

For a compact connected space satisfying Theorem~\ref{maximum}, take a continuous PC reference $\gamma$ and a continuous obstacle $u$. Since $\gamma$ is constant, the envelope over \emph{continuous} PC minorants is the constant
\[\min_X(\gamma+u).
\]
If a finite bounded rational polyhedral cell decomposition compatible with the local charts is given and $\gamma,u$ have literal integral affine slopes and rational constants on a finite rational subdivision, their minimum value is rational: it is attained at a rational vertex of a cell of a common bounded subdivision. The constant envelope is consequently literal PA and has a constant decreasing regularizing sequence. In positive dimension its raw iterated-divisor measure is zero, because intersection with a constant is zero. Its contact integral is zero. These are direct finite-PA assertions and introduce no operator extension. An envelope allowing competitors without an explicit continuity hypothesis requires a separate identification of those competitors with continuous PC functions; that identification is not granted here.

In dimension zero a finite globally embedded space is a finite set of rational points. Every PCC is trivial on its singleton face. A literal PA reference and obstacle therefore have exact envelope $\gamma+u$ on each point, the raw zero-fold measure is the given weight cycle, and the contact gap is identically zero.

Theorems~\ref{noncompact} and \ref{maximum} clarify compact-category boundaries. They do not prove real-PA regularity on noncompact singular spaces, the global realization of local noncontact germs, or a necessary-and-sufficient higher-dimensional contact criterion. In particular they do not close original P05.1 or P05.2. The positive-dimensional real sector remains an original target.

\appendix
\section{Source application and reproducibility}
The sole cited primary version is \cite{BMPS}. Definitions 2.1--2.4 were read in the literal PDF text, lines 222--265. Definitions 2.18--2.27 and the full local-star Lemma 2.24 proof were read in lines 443--527; Definitions 3.2--3.3 were read in lines 880--902. The fixed source requires finite globally embedded complexes and strictly positive top cycles. Its toric compactification is a different space in additional strata. These definitions are paraphrased above; every needed geometric assertion is proved independently. No proof of its Monge--Amp\`ere extension, its orthogonality theorem, or the cited automatic-continuity result was imported. All auxiliary Euclidean choices disappear from the geometric conclusions.

This is an unreviewed supporting manuscript. Compilation and provenance hashes do not certify its mathematics. The exact saved source can be compiled twice using an existing \texttt{pdflatex}; no numerical experiment is needed for the divergence identities or Jensen argument.

\begin{thebibliography}{1}
\bibitem{BMPS} A. M. Botero, E. Mazzon, and L. Pille--Schneider,
\emph{Monge--Amp\`ere measures on balanced polyhedral spaces},
arXiv:2603.08664v1, submitted 9 March 2026.
\url{https://arxiv.org/pdf/2603.08664v1}.
\end{thebibliography}
\end{document}
